% Source: https://artofproblemsolving.com/community/c3h1917408_tricky_parabola_questions
% Stable identity: aops-topic-1917408-q02
The parabolas meet where $x^2-b=-x^2+b$, i.e.\ $x^2=b$, at $P(\sqrt b,0)$ and $Q(-\sqrt b,0)$. Their vertices are $(0,-b)$ and $(0,b)$.

The quadrilateral with these four vertices has diagonals along the axes, of lengths $2\sqrt b$ and $2b$; they are perpendicular and bisect each other at the origin, so it is a rhombus (drawn below for $b=9$):

\begin{asy}
usepackage("amsmath");
size(5.5cm);
defaultpen(fontsize(10pt));
import graph;
real g1(real x){return x^2-9;} real g2(real x){return 9-x^2;}
fill((3,0)--(0,9)--(-3,0)--(0,-9)--cycle,palegreen);
draw((-5.5,0)--(5.5,0),Arrow(6)); draw((0,-10.5)--(0,11),Arrow(6));
label("$x$",(5.5,0),E); label("$y$",(0,11),N);
draw(graph(g1,-4.1,4.1),blue+linewidth(1)); draw(graph(g2,-4.1,4.1),red+linewidth(1));
draw((3,0)--(0,9)--(-3,0)--(0,-9)--cycle,linewidth(0.8));
dot((3,0));dot((-3,0));dot((0,9));dot((0,-9));
label("$P$",(4.45,0.9)); label("$Q$",(-4.45,0.9));
label("$(0,b)$",(0,9),NE); label("$(0,-b)$",(0,-9),SE);
label("$y=x^2-b$",(4.1,7.81),E,blue); label("$y=b-x^2$",(4.1,-7.81),E,red);
label("\mbox{}",(min(currentpicture,user=true)+max(currentpicture,user=true))/2,black); // keeps the text after the figure black when LaTeX typesets the labels inline
\end{asy}

Its area is half the product of the diagonals:
\[
\dfrac{1}{2}\cdot2\sqrt b\cdot2b=2b\sqrt b=54\iff b\sqrt b=27\iff b^{\dfrac{3}{2}}=27\iff b=\boxed{9}.
\]
